\section{Resumo}
O presente trabalho experimental teve como principais objetivos a calibração de uma sonda de efeito de Hall, o estudo do campo magnético gerado por bobinas de Helmholtz e a verificação experimental do princípio da sobreposição de campos magnéticos.
A metodologia consistiu na medição da tensão de Hall para diferentes valores de campo magnético conhecido, seguida de regressão linear para obtenção da curva de calibração, e na posterior determinação experimental do campo ao longo do eixo das bobinas em várias configurações (correntes no mesmo sentido e em sentidos opostos).

A calibração revelou uma relação linear $V_H =\Slope,B$, com coeficiente de determinação $R^2 =\Rsquare$, e desvio padrão de $\StdDev\%$ e o termo independente compatível com zero, que evidenciam uma elevada precisão e boa exatidão.